\section{A symmetric hierarchy of nearly first eigenfunctions}
\label{sec:hierarchy}

The prediction curve from Section~\ref{sec:prediction} determines much more
than the norm of a conditional mean.  Its Gaussian Bessel inequality is
asymptotically saturated at every fixed derivative order.  We use this
saturation to construct symmetric tensors whose derivatives remain nearly
first eigenfunctions.  We then show that every tensor has negligible energy
in the curvature directions of $\mu_s$.  Throughout this section, coefficient
norms are unnormalized Euclidean norms, including for tensors of high order.

For a tensor field with $d$ slots, differentiation appends a last slot.  A
matrix subscript, such as $K_1$, specifies the slot on which it acts.  Write
$\|F\|_2^2=\E_\eta|F|^2$ and
$\langle F,G\rangle_2=\E_\eta\langle F,G\rangle$.  Every occurrence of $A$ on a
tensor means its componentwise action.  The following formulation records
precisely the inputs used in the construction; the dimension and $I$ may
vary along the sequence.

\begin{proposition}[Symmetric tensor hierarchy]\label{prop:hierarchy}
Consider a sequence of smooth positive probability densities $\eta$ on
finite-dimensional Euclidean spaces, with finite second moments.  Let
$A=-\Delta-D\log\eta\cdot D$, and suppose that $\eta$ satisfies Poincar\'e
with constant at most $\lambda^{-1}$, where $\lambda\to1$.  Fix $s=1/8$.
Suppose smooth vector fields $b,w$ and a symmetric matrix field $K$ satisfy
\[
 D\log\eta=-\frac{y-b(y)}s+\omega(y),\qquad
 Db=\Id-sK,\qquad 0\le K\le s^{-1}\Id,
\]
where either $\omega=w$ throughout the sequence or $\omega=0$ throughout
the sequence.  Assume $b,w\in L^2(\eta)$, $w$ is a gradient, and, with
$I=\|w\|_2^2>0$,
\begin{equation}\label{eq:hierarchy-initial}
 w\in\Dom A,\qquad
 \|(A-1)w\|_2=o(\sqrt I),\qquad
 |\E_\eta w|=o(\sqrt I),\qquad
 \langle w,Kw\rangle_2=o(I).
\end{equation}
Assume also the following prediction and projection properties.
\begin{enumerate}[label=\textup{(\roman*)}]
\item For $Y\sim\eta$ and an independent standard Gaussian $G$, set
\[
 U_r=Y+\sqrt r\,G,\quad h_r(u)=\E[w(Y)\mid U_r=u],\quad
 m(r)=I^{-1}\E|h_r(U_r)|^2.
\]
For every fixed $0<r<1/4$, one has $m(r)\to e^{-r}$.
\item For every fixed $c<\infty$,
\begin{equation}\label{eq:hierarchy-input-projection}
 \sup_{\substack{\Pi^2=\Pi=\Pi^*\ \mathrm{deterministic}\\
                   \rank\Pi\le cI}}
       \|\Pi w\|_2^2=o(I).
\end{equation}
For every $\epsilon>0$ there is a deterministic projection $\Pi_\epsilon$
of rank at most $C_\epsilon I$ such that
\begin{equation}\label{eq:hierarchy-input-outliers}
 \left\|\E[(1-\Pi_\epsilon)bb^*(1-\Pi_\epsilon)]\right\|_{\op}
 +\left\|\E[(1-\Pi_\epsilon)\omega\omega^*
                  (1-\Pi_\epsilon)]\right\|_{\op}\le\epsilon
\end{equation}
eventually along the sequence.
\item If $\omega=w$, the conditional covariance
$C_r=\Cov(w(Y)\mid U_r)$ satisfies, for every fixed $r\in(0,1/4)$ and
$a>0$,
\begin{equation}\label{eq:hierarchy-input-tail}
 \E\tr\bigl(C_r\1_{\{C_r>a\}}\bigr)=o(I).
\end{equation}
No conditional covariance-tail assumption is made when $\omega=0$.
\end{enumerate}
Then there are totally symmetric tensor fields $T_l$ of order $l+1$, for
every fixed integer $l\ge0$, with $T_0=w$, such that
\begin{equation}\label{eq:hierarchy-eigen}
 \begin{gathered}
 \|T_l\|_2^2=(1+o(1))I,\qquad
 \|(A-1)T_l\|_2=o(\sqrt I),\\
 \|DT_l-T_{l+1}\|_2=o(\sqrt I).
 \end{gathered}
\end{equation}
\begin{equation}\label{eq:hierarchy-curvature}
 \langle T_l,K_1T_l\rangle_2=o(I).
\end{equation}
\begin{equation}\label{eq:hierarchy-projection}
 \sup_{\substack{\Pi^2=\Pi=\Pi^*\ \mathrm{deterministic}\\
                   \rank\Pi\le cI}}
      \|\Pi_1T_l\|_2^2=o(I)
 \qquad\text{for every fixed }c<\infty.
\end{equation}
Every $T_l$ belongs componentwise to $\Dom A$.  For $l\ge1$ it can be
chosen in $\bigcap_{p\ge1}\Dom A^p$.  By symmetry, the specified slot in
Equations~\eqref{eq:hierarchy-curvature} and
\eqref{eq:hierarchy-projection} may be any slot.  All assertions can be
made simultaneous for any fixed finite collection of levels and parameters.
\end{proposition}

In the application, take $b=b_s$ and $K=K_s$.  Proposition~\ref{prop:critical-sequence}
supplies the drift, Poincar\'e, and initial estimates;
Proposition~\ref{prop:prediction} and Lemma~\ref{lem:posterior-outliers}
supply the remaining inputs.  In the gap branch, $\omega=0$ and
$\|\E bb^*\|_{\op}=o(1)$, so one may take $\Pi_\epsilon=0$ in
Equation~\eqref{eq:hierarchy-input-outliers}.  The normalization $I\to0$
makes Equation~\eqref{eq:hierarchy-input-projection} automatic: eventually
the only projection of rank at most $cI$ is zero.  Thus no strict-branch
probe estimate has been added to the gap-branch hypotheses.

We prove the proposition in three steps.  First we construct the
eigenfunction and derivative hierarchy.  Second we expand its evolution
along stationary diffusion and, in the strict branch, deduce operator-norm
tail estimates.  Finally those estimates justify a curvature integration
by parts.

\subsection{Gaussian Bessel saturation and the induction}

We use the classical correspondence between Hermite polynomials and
multiple Wiener integrals \cite[Sections~3--5]{Ito1951}, with the
ordered-index normalization recorded below.  The saturation argument
will come from the prediction curve, in addition to this orthogonality.

For a standard Gaussian vector $G\in\R^n$, define its symmetric Hermite
tensors $\mathfrak h_j(G)$ by
\begin{equation}\label{eq:hierarchy-hermite-definition}
 e^{\xi\cdot G-|\xi|^2/2}
 =\sum_{j\ge0}\frac1{j!}\,\mathfrak h_j(G):\xi^{\otimes j}.
\end{equation}
Here a colon contracts all indicated slots and $\mathfrak h_0=1$.
Comparing coefficients in the expectation of the product of two such
exponentials gives, for symmetric order-$j$ and order-$k$ tensors $a,b$,
\begin{equation}\label{eq:hierarchy-hermite-isometry}
 \E[(a:\mathfrak h_j(G))(b:\mathfrak h_k(G))]
 =\1_{\{j=k\}}j!\langle a,b\rangle.
\end{equation}
This identity uses ordered coefficient indices; it explains all factorials
below.

\begin{lemma}[Bessel inequality for the prediction]\label{lem:hierarchy-bessel}
For the data of Proposition~\ref{prop:hierarchy}, define
$g_r(y)=\E_G h_r(y+\sqrt r\,G)$.  For each $r>0$,
\begin{equation}\label{eq:hierarchy-bessel}
 \langle w,g_r\rangle_2=m(r)I,\qquad
 \sum_{j\ge0}\frac{r^j}{j!}\|D^jg_r\|_2^2\le m(r)I.
\end{equation}
In particular, every fixed derivative appearing here is in the form
domain of $A$.
\end{lemma}

\begin{proof}
The first equality is conditional expectation:
\[
 \E\langle w(Y),h_r(U_r)\rangle=\E|h_r(U_r)|^2.
\]
For fixed $y$, Gaussian integration by parts identifies the order-$j$
Hermite coefficient of $h_r(y+\sqrt r\,G)$ as
$r^{j/2}D^jg_r(y)$.  Its orthogonal projection onto that chaos is
\[
 \frac{r^{j/2}}{j!}D^jg_r(y):\mathfrak h_j(G),
\]
where the vector output slot is not contracted.  By
Equation~\eqref{eq:hierarchy-hermite-isometry}, the squared norm of this
projection is $r^j|D^jg_r(y)|^2/j!$.  Bessel's inequality, integrated in
$y\sim\eta$, proves the asserted sum.

These calculations require only
$\E|h_r(Y+\sqrt rG)|^2<\infty$.  To justify the derivatives, on a compact
set of starting points use positivity of $\eta$ on a slightly larger
compact set.  Integrating translated Gaussian kernels over that larger
set bounds each fixed differentiated kernel by a constant times the
mixture density of $U_r$.  The enlargement absorbs the polynomial factors
from differentiation.  Thus the kernel integrals are classically smooth
locally; alternatively all identities first hold distributionally.  The
Bessel bounds for orders $j$ and $j+1$ put $D^jg_r$ and its weak derivative
in $L^2(\eta)$, hence in the form domain by cutoff and local smoothing.
The local constants need not be uniform in the dimension: only the exact
integrated Bessel inequality is used quantitatively.
\end{proof}

Two elementary Hilbert-space facts will prevent dimension factors and
unjustified differentiations in the induction.

\begin{lemma}[Cross-covariance and spectral smoothing]
\label{lem:hierarchy-smoothing}
If $F$ is a square-integrable random variable with values in a
finite-dimensional Hilbert space and $Z$ is a
centered Euclidean random vector with $\Cov Z\le c\Id$, then
\begin{equation}\label{eq:hierarchy-crosscov}
 |\E(F\otimes Z)|_{\HS}^2\le c\E|F|^2.
\end{equation}
Next, let $A$ be a sequence of nonnegative Dirichlet operators with kernel
the constants and spectral gap at least $\lambda\to1$.  If fields $S$ in
their form domains satisfy
\[
 |\E S|^2=o(I),\qquad \|S\|_2^2=(1+o(1))I,\qquad
 \|DS\|_2^2\le(1+o(1))I,
\]
then $\widetilde S=e^{1-A}S$ satisfies
\begin{equation}\label{eq:hierarchy-smoothing}
 \|\widetilde S-S\|_2+\|D(\widetilde S-S)\|_2=o(\sqrt I),
 \qquad \|(A-1)\widetilde S\|_2=o(\sqrt I).
\end{equation}
\end{lemma}

\begin{proof}
The map $u\mapsto u\cdot Z$ from Euclidean space to scalar $L^2$ has
operator norm at most $\sqrt c$.  Apply its adjoint to each scalar
coefficient of $F$ and sum the squared bounds.  This proves
Equation~\eqref{eq:hierarchy-crosscov} without a factor equal to the number
of coordinates of either variable.

For the second assertion, sum the coefficient spectral measures of
$S-\E S$ and call the resulting measure $\nu$.  Then
\[
 \int(x-\lambda)\,d\nu(x)
 =\|DS\|_2^2-\lambda\|S-\E S\|_2^2=o(I),
\]
where the left side is nonnegative.  For every fixed $\delta>0$, the
spectral mass and the spectral energy outside
$[1-\delta,1+\delta]$ are $o(I)$: the part below $1-\delta$ is absent
eventually, and on $x\ge1+\delta$ both $1$ and $x$ are bounded by a
$\delta$-dependent constant times $x-\lambda$.  The constant component
has squared norm $o(I)$ as well.

On the remaining interval the multipliers $e^{1-x}-1$ and
$(x-1)e^{1-x}$ tend uniformly to zero as $\delta\downarrow0$.  Outside
it, $(1+x)|e^{1-x}-1|^2\le C(1+x)$ and
$(x-1)^2e^{2-2x}\le C$.  Integrate these inequalities against the
spectral measure, first take the sequence limit and then let
$\delta\downarrow0$.  This gives Equation~\eqref{eq:hierarchy-smoothing}.
The same functional calculus shows that $e^{1-A}S$ belongs to every
$\Dom A^p$.
\end{proof}

\begin{lemma}[Construction before the curvature estimate]
\label{lem:hierarchy-construction}
Under the hypotheses of Proposition~\ref{prop:hierarchy}, there exist
totally symmetric fields $T_l$ satisfying
Equations~\eqref{eq:hierarchy-eigen} and
\eqref{eq:hierarchy-projection}, with the stated operator domains.
Moreover $|\E T_l|=o(\sqrt I)$ for every fixed $l$.
\end{lemma}

\begin{proof}
Start with $T_0=w$.  Suppose the assertions have been constructed through
level $l$, and put $L=l+1$ and $V=DT_l$.  Self-adjointness gives
\begin{equation}\label{eq:hierarchy-raw-derivative}
 \|V\|_2^2=\langle T_l,AT_l\rangle_2=(1+o(1))I.
\end{equation}
If $\Pi$ is an admissible deterministic projection acting on an existing
tensor slot, it commutes with $A$ and differentiation.  Consequently
\[
 \|\Pi_1V\|_2^2
 =\langle\Pi_1T_l,A\Pi_1T_l\rangle_2
 \le\|\Pi_1T_l\|_2^2+
       \|\Pi_1T_l\|_2\|(A-1)T_l\|_2=o(I).
\]
The error is uniform over projections of rank at most $cI$ for fixed $c$.

Fix $0<r<1/4$.  Set $V_j=T_j$ for $0\le j\le l$ and $V_L=V$.
Successive form integration by parts gives
\begin{equation}\label{eq:hierarchy-correlations}
 I^{-1}\langle D^jg_r,V_j\rangle_2\longrightarrow e^{-r}
 \qquad(0\le j\le L).
\end{equation}
Indeed the assertion at $j=0$ is
Equation~\eqref{eq:hierarchy-bessel}.  For $j\le l$,
\[
 \langle D^{j+1}g_r,DT_j\rangle_2
 =\langle D^jg_r,AT_j\rangle_2
 =\langle D^jg_r,T_j\rangle_2+o(I).
\]
Replacing $DT_j$ by $T_{j+1}$ when $j<l$ costs $o(I)$.
All test derivatives and their form energies are finite by
Lemma~\ref{lem:hierarchy-bessel}; its bound controls every error for
fixed $r$ and $j$.  No uniform estimate as $r\downarrow0$ has been used.

Here is the precise Bessel budget.  Equation~\eqref{eq:hierarchy-correlations}
and $\|V_j\|_2^2/I\to1$ imply
\[
 \liminf I^{-1}e^{2r}\|D^jg_r\|_2^2\ge1
 \qquad(0\le j\le L).
\]
Multiply Equation~\eqref{eq:hierarchy-bessel} by $e^{2r}/I$ and subtract
these minimum contributions.  The resulting bounds are
\begin{align}
 \limsup\frac{\|e^rD^Lg_r-V\|_2^2}{I}
 &\le\frac{L!}{r^L}
       \left(e^r-\sum_{j=0}^{L}\frac{r^j}{j!}\right)=O_L(r),
 \label{eq:hierarchy-budget-error}\\
 \limsup\frac{e^{2r}\|D^{L+1}g_r\|_2^2}{I}
 &\le\frac{(L+1)!}{r^{L+1}}
       \left(e^r-\sum_{j=0}^{L}\frac{r^j}{j!}\right)=1+O_L(r).
 \label{eq:hierarchy-budget-energy}
\end{align}
For the first line, expand the squared difference and use the correlation
at order $L$; its normalized limit subtracts exactly one from the
normalized squared derivative norm.  Thus the exponential-series tail
accounts for the entire excess, rather than merely bounding the norms.

Choose $r\downarrow0$ sufficiently slowly along the sequence in these
two estimates.  The field $S_0=e^rD^Lg_r$ then satisfies
\[
 \|S_0-V\|_2=o(\sqrt I),\qquad
 \|DS_0\|_2^2\le(1+o(1))I.
\]
It is symmetric in its last $L$ slots; $V$ is symmetric in its first
$L$ slots.  For $L\ge2$, the permutations of those two sets of slots
generate every permutation of the $L+1$ slots.  Indeed their adjacent
transpositions together are all the adjacent transpositions.  If a
permutation fixes $S_0$, its action changes $V$ by at most
$2\|S_0-V\|_2$; a permutation fixing $V$ costs zero.  Expressing each
permutation as a fixed finite word therefore shows that $V$ is
$o(\sqrt I)$ from its full symmetrization.  For $L=1$, $V=Dw$ is
already symmetric because $w$ is a gradient.  Let $S$ be the full
symmetrization of $S_0$.  Averaging coefficient permutations is an
orthogonal contraction and commutes with differentiation, so
\begin{equation}\label{eq:hierarchy-symmetric-approximant}
 \|S-V\|_2=o(\sqrt I),\qquad
 \|S\|_2^2=(1+o(1))I,\qquad
 \|DS\|_2^2\le(1+o(1))I.
\end{equation}
The preceding projection estimate for $V$ now holds on its last slot
too, since a coefficient-slot permutation moves it to an old slot with
an $o(\sqrt I)$ error uniform over projections.

To apply Poincar\'e sharply we still need to center $V$.  Write
$C_y=\Cov(T_l,y)$, $C_b=\Cov(T_l,b)$, and
$C_\omega=\Cov(T_l,\omega)$, with the vector variable occupying the
last slot.  Testing the equation for $AT_l$ against centered coordinates
gives
\begin{equation}\label{eq:hierarchy-two-center-identities}
 \E V=C_y+E_l,\qquad |E_l|=o(\sqrt I),\qquad
 \E V=s^{-1}(C_y-C_b)-C_\omega.
\end{equation}
For clarity, the first exact identity before the error is
$\E V=\E[(AT_l)\otimes(y-\E y)]$.  Poincar\'e gives
$\Cov_\eta y\le\lambda^{-1}\Id$, and
Equation~\eqref{eq:hierarchy-crosscov} bounds its error by
$\lambda^{-1/2}\|(A-1)T_l\|_2$.  This is the needed dimension-free
estimate.  The second identity follows by integration by parts using
$Ay=(y-b)/s-\omega$.  These coordinate tests are legitimate in the
form domain; their distributional $A$-images are in $L^2$ by the stated
second-moment assumptions, hence they are also in $\Dom A$.

Project the last slot onto $1-\Pi_\epsilon$.  The two covariance terms
involving $b$ and $\omega$ have norm at most $C\sqrt{\epsilon I}$ by
Equations~\eqref{eq:hierarchy-input-outliers} and
\eqref{eq:hierarchy-crosscov}.  Subtracting the two identities in
Equation~\eqref{eq:hierarchy-two-center-identities}, and using $s\ne1$,
therefore bounds this part of $\E V$ by
$C_s\sqrt{\epsilon I}+o(\sqrt I)$.  On the $\Pi_\epsilon$ part,
\[
 |(\Pi_\epsilon)_{L+1}\E V|
 \le\|(\Pi_\epsilon)_{L+1}V\|_2=o(\sqrt I)
\]
by the projected-energy estimate and $\rank\Pi_\epsilon\le C_\epsilon I$.
First take the sequence limit and then let $\epsilon\downarrow0$.
This proves $|\E V|=o(\sqrt I)$, and thus $|\E S|=o(\sqrt I)$.

Lemma~\ref{lem:hierarchy-smoothing} applies to $S$.  Define
\[
 T_{l+1}=e^{1-A}S.
\]
It is symmetric, lies in every operator-power domain, is
$o(\sqrt I)$ from $V$, and satisfies the strong near-eigenfunction
estimate.  Its mean is $e\E S=o(\sqrt I)$.  Finally, coefficient
projections commute with the semigroup, and $\|e^{1-A}\|_{2\to2}\le e$.
Together with $S-V=o_{L^2}(\sqrt I)$, this preserves the uniform
projected-energy estimates.  The induction is complete.  Each step used
only a finite collection of limits; diagonal choices give the assertion
for all fixed levels.
\end{proof}

\subsection{Expansion along stationary diffusion}

The fields just constructed retain the same norm under repeated
differentiation.  This makes their stochastic Taylor expansion summable
with a factorial remainder.  We use the ordered-integral form of the
Wiener--It\^o correspondence \cite[Section~5]{Ito1951}.  Let $Y_t$ be
the stationary diffusion with
invariant density $\eta$,
\[
 dY_t=\sqrt2\,dB_t+D\log\eta(Y_t)\,dt,\qquad Y_0\sim\eta,
\]
where $Y_0$ is independent of the driving Brownian motion.

\begin{lemma}[Hermite expansion of the stationary evolution]
\label{lem:hermite-expansion}
Let symmetric fields $T_l$ satisfy Equation~\eqref{eq:hierarchy-eigen}.
For fixed integers $l,m\ge0$ and fixed $t>0$, set $N=B_t/\sqrt t$.  Then
\begin{equation}\label{eq:hierarchy-diffusion-expansion}
 T_l(Y_t)=e^{-t}\sum_{j=0}^m\frac{(2t)^{j/2}}{j!}
       T_{l+j}(Y_0):\mathfrak h_j(N)+\mathcal E_{l,m,t},
\end{equation}
where the colon contracts the last $j$ tensor slots and
\begin{equation}\label{eq:hierarchy-diffusion-remainder}
 \limsup\frac{\|\mathcal E_{l,m,t}\|_2}{\sqrt I}
 \le\frac{(2t)^{(m+1)/2}}{\sqrt{(m+1)!}}.
\end{equation}
For $0\le t\le T<\infty$ the same bound holds uniformly with $t$
replaced on the right by $T$ and with the expansion at $t=0$ interpreted
by continuity.  The limits are also simultaneous for any fixed finite
set of starting levels.  In particular the normalized remainder tends
to zero when the sequence limit is followed by $m\to\infty$.
\end{lemma}

\begin{proof}
The operator-domain It\^o formula from Lemma~\ref{lem:dirichlet-core}
gives, for each fixed $d$,
\begin{align*}
 e^vT_d(Y_v)
 ={}&T_d(Y_0)+\sqrt2\int_0^v e^uT_{d+1}(Y_u):dB_u\\
 &+\int_0^v e^u(1-A)T_d(Y_u)\,du
   +\sqrt2\int_0^v e^u(DT_d-T_{d+1})(Y_u):dB_u.
\end{align*}
Stationarity, Cauchy--Schwarz for the time integral, and It\^o isometry
show that the last two terms are $o(\sqrt I)$, uniformly for
$v\in[0,T]$ in squared-mean norm.  Substitute this formula for the
integrand at the next level, and repeat $m+1$ times.  For fixed $l,m,T$,
every replacement error remains $o(\sqrt I)$; iterated isometry sums
coefficient squares and introduces no dimension factor.

After multiplication by $e^{-t}$, the remaining integral has coefficient
$2^{(m+1)/2}$, is over the ordered time simplex of length $t$, and has
innermost integrand $e^{u-t}T_{l+m+1}(Y_u)$.  Since $u\le t$,
its squared norm is at most
\[
 2^{m+1}\frac{t^{m+1}}{(m+1)!}
       \|T_{l+m+1}\|_2^2
 =(1+o(1))I\frac{(2t)^{m+1}}{(m+1)!}.
\]
This proves the remainder estimate and its uniform version.
The other terms have coefficients evaluated at $Y_0$.  Symmetry in
their Brownian slots turns their ordered iterated integrals into
$t^{j/2}\mathfrak h_j(N)/j!$.  To verify the factor directly, expand
the Brownian exponential martingale
$\exp(\xi\cdot B_t-t|\xi|^2/2)$ by iterated It\^o integration and
compare its coefficients with
Equation~\eqref{eq:hierarchy-hermite-definition}.  The additional
$\sqrt2$ at each stochastic integration gives exactly
$(2t)^{j/2}/j!$ in
Equation~\eqref{eq:hierarchy-diffusion-expansion}.
\end{proof}

\subsection{Operator tails in the strict branch}

For a tensor $F$, a one-slot flattening is the matrix whose row index is
one specified tensor slot and whose column index is the tuple of all
remaining slots.  Its Hilbert--Schmidt norm is exactly $|F|$.

In the strict branch, put $Q=Dw$. The drift decomposition gives
$D^2(-\log\eta)=K-Q$. Viewing $DF$ as a matrix with its derivative
slot as row index, commuting a derivative with $A$ formally gives
\[
 D(AF)=A(DF)+(K-Q)DF.
\]
For $F=T_l$, both $Q$ and $DF$ belong to $L^2$, but their product
$QDF$ need not belong to $L^2$. We cannot yet use this formula to
deduce the curvature estimate. The conditional covariance bound will
instead let us approximate $DF$ by tests $P$ whose uniform operator
norms tend to zero. Then
$\|QP\|_2\le\|P\|_{\op,\infty}\|Q\|_2=o(\sqrt I)$,
which is the product bound needed for the justified identity below.
This step is needed only when $\omega=w$; in the gap branch $Q=0$.

\begin{lemma}[Vanishing operator tails]\label{lem:hierarchy-operator-tails}
In the strict branch of Proposition~\ref{prop:hierarchy}, $Q=Dw$ and
the fields $T_l$ constructed in Lemma~\ref{lem:hierarchy-construction}
satisfy the following property.  For every fixed $\delta>0$, each
one-slot flattening $F$ of $T_l$, $l\ge1$, and also $F=Q$, obeys
\begin{equation}\label{eq:hierarchy-operator-tail}
 \E\tr\bigl(FF^*\1_{\{FF^*>\delta^2\}}\bigr)=o(I).
\end{equation}
Equivalently, clipping every singular value at any fixed positive cutoff
changes the field by $o(\sqrt I)$ in $L^2$.
\end{lemma}

\begin{proof}
We first prove the assertion for $Q$.  Couple the heat observations by
$U_v=Y_0+W_v$, where $W$ is a standard Brownian motion independent of
$Y_0\sim\eta$, and retain $h_v=\E[w(Y_0)\mid U_v]$.  Put
$H_v=Dh_v(U_v)$ wherever this derivative is square integrable.  The
conditional covariance has the exact representation
\begin{equation}\label{eq:hierarchy-conditional-bracket}
 C_r=\int_0^r\E[H_vH_v^*\mid U_r]\,dv,
 \qquad
 \int_0^r\E|H_v|^2\,dv=(1-m(r))I.
\end{equation}
To check the first identity and its normalization, let $p_v$ denote the
density of $U_v$.  The entries of $p_vh_v$ solve the heat equation with
generator $\Delta/2$.  Direct differentiation gives the matrix identity
\[
 (\partial_v-\tfrac12\Delta)(p_vh_vh_v^*)
 =-p_vDh_v(Dh_v)^*.
\]
Integrate it against the later heat kernel and divide by $p_r$.  The
initial term becomes $\E[ww^*\mid U_r]$, and the final term is
$h_rh_r^*$, which proves the first formula.  Taking the trace and
expectation proves the second formula.  For nonsmooth or unbounded
initial $w$, perform the calculation on scalar projections with convex
truncations increasing to the square, and then use polarization.  Spatial
cutoffs and the square-integrability of $w$ justify the positive-time
passage.  At time zero, approximation of $w$ in $L^2(\eta)$ by bounded
continuous functions gives
$\E|w-\E[w\mid U_v]|^2\to0$.  Thus there is no missing boundary term
in Equation~\eqref{eq:hierarchy-conditional-bracket}.

By Equation~\eqref{eq:hierarchy-initial},
$\|Q\|_2^2=\langle w,Aw\rangle_2=(1+o(1))I$.
For almost every $v>0$, convolution differentiation and form integration
by parts give
\begin{equation}\label{eq:hierarchy-derivative-correlation}
 \E\langle Q(Y_0),H_v\rangle
 =\langle Dw,Dg_v\rangle_2
 =\langle Aw,g_v\rangle_2=m(v)I+o(I).
\end{equation}
The error is uniform in $v$: Jensen's inequality gives
$\|g_v\|_2\le\sqrt I$, so it is bounded by
$\|(A-1)w\|_2\sqrt I$.  The convolution derivative identity follows
first with cutoffs and bounded functions, then by the square-integrable
derivative bounds in Equation~\eqref{eq:hierarchy-conditional-bracket}.

Fix a small $z>0$ with $2z<1/4$.  Combining the last two equations and
using $0\le m(v)\le1$ yields the exact limiting discrepancy
\begin{align}
 \lim\frac1I\int_0^{2z}\E|Q(Y_0)-H_v|^2\,dv
 &=2z+1-e^{-2z}-2\int_0^{2z}e^{-v}\,dv\\
 &=2z-1+e^{-2z}=O(z^2).
 \label{eq:hierarchy-integrated-discrepancy}
\end{align}
Dominated convergence is applicable to $m(v)$, so this calculation does
not require uniform convergence of the prediction curve near zero.

We select a time in $[z,2z]$ while preserving a covariance-tail bound.
For a fixed $\kappa>0$, define
\[
 a(r)=I^{-1}\E\tr\left(\frac{C_r}{r}
                   \1_{\{C_r/r>\kappa\}}\right).
\]
For fixed $r\in[z,2z]$, Equation~\eqref{eq:hierarchy-input-tail} gives
$a(r)\to0$, and $a(r)\le1/z$ follows from $\E\tr C_r\le I$.
Consequently $\int_z^{2z}a(r)\,dr\to0$.  On the other hand,
Equation~\eqref{eq:hierarchy-integrated-discrepancy} and Markov's
inequality show that the set of times satisfying
\[
 \E|Q(Y_0)-H_r|^2\le C zI
\]
has measure at least $z/2$, with an absolute $C$ and for sufficiently
far terms of the sequence.  The set where $a(r)$ exceeds a quantity
tending to zero has measure tending to zero.  Select $r\in[z,2z]$
in the intersection.  Thus $r$ may vary with the sequence index, but
\begin{equation}\label{eq:hierarchy-selected-time}
 \E|Q(Y_0)-H_r|^2\le CzI,\qquad
 \E\tr\left(\frac{C_r}{r}\1_{\{C_r/r>\kappa\}}\right)=o(I).
\end{equation}
All constants in this selection are for fixed $z,\kappa$.

We compare $B=H_rH_r^*$ with $C=C_r/r$ in trace norm.  For matrices
$U,V$,
\begin{equation}\label{eq:hierarchy-product-trace}
 \|UU^*-VV^*\|_1
 \le |U-V|_{\HS}(|U|_{\HS}+|V|_{\HS}).
\end{equation}
Compare both $B$ and $C$ with $\E[QQ^*\mid U_r]$.  For $B$, conditional
Jensen's inequality and the first bound in
Equation~\eqref{eq:hierarchy-selected-time} give an $O(\sqrt z)I$
bound on the expected trace-norm difference.  For $C$, use the first
identity in Equation~\eqref{eq:hierarchy-conditional-bracket}, followed
by Equation~\eqref{eq:hierarchy-product-trace} and Cauchy--Schwarz in
both probability and time.  The integrated discrepancy is $O(z^2)I$,
whereas
\[
 \int_0^r\E(|H_v|+|Q|)^2\,dv
 \le 2(1-m(2z))I+4z\|Q\|_2^2.
\]
Here we used $r\le2z$ and the nonnegative integral in the second
identity of Equation~\eqref{eq:hierarchy-conditional-bracket}.
The normalized limsup of the right side is $O(z)$, by prediction at
the fixed time $2z$. No prediction limit at the selected varying time
is needed.
Division by $r\ge z$ therefore also costs at most $O(\sqrt z)I$.
We have proved, in normalized limsup,
\begin{equation}\label{eq:hierarchy-bracket-comparison}
 \E\|B-C\|_1\le O(\sqrt z)I.
\end{equation}

Fix $\delta>0$, take $\kappa=\delta^2/2$ above, and set
$P=\1_{\{B>\delta^2\}}$.  Positivity of $C$ implies
\[
 \tr(BP)\le\frac{\delta^2}2\rank P
      +\tr\bigl(C\1_{\{C>\delta^2/2\}}\bigr)+\|B-C\|_1.
\]
Since $\delta^2\rank P\le\tr(BP)$, the first term is absorbed into
the left side.  Equations~\eqref{eq:hierarchy-selected-time} and
\eqref{eq:hierarchy-bracket-comparison} bound the remaining high
singular-value energy by $O(\sqrt z)I+o(I)$.  Clip $H_r$ at $\delta$;
its squared distance from the operator-norm ball of radius $\delta$ is
at most this energy.  The first estimate in
Equation~\eqref{eq:hierarchy-selected-time} then gives the same bound,
up to constants, for the distance of $Q$ from that ball.  Let $z\downarrow0$
after the sequence limit.  The squared distance is $o(I)$ for every
fixed $\delta>0$.  Applying this distance assertion at cutoff
$\delta/2$ bounds the energy of singular values above $\delta$ by
four times the squared distance, proving
Equation~\eqref{eq:hierarchy-operator-tail} for $Q$.  Since
$\|T_1-Q\|_2=o(\sqrt I)$, it also proves the assertion for $T_1$.

It remains to propagate this estimate to higher levels.  Fix $l\ge1$
and $t>0$.  Apply Lemma~\ref{lem:hermite-expansion} with starting level
$1$.  Conditional on $Y_0$, the Hermite chaoses are orthogonal even
after applying an arbitrary measurable projection $P(Y_0)$ to the
first tensor slot.  The chaos containing $T_l(Y_0)$ has squared
coefficient
\[
 c_{l,t}=e^{-2t}\frac{(2t)^{l-1}}{(l-1)!}>0.
\]
Use the expansion to order $m\ge l-1$, first take the sequence limit,
and then let $m\to\infty$ in its remainder bound.  The resulting
inequality is
\begin{equation}\label{eq:hierarchy-random-projection}
 \liminf\frac{
   \E|P(Y_0)_1T_1(Y_t)|^2-c_{l,t}\E|P(Y_0)_1T_l(Y_0)|^2}{I}\ge0.
\end{equation}
The remainder estimates are uniform over such projections, since their
operator norms are at most one.

Choose $P(Y_0)$ to project onto the singular values greater than a fixed
$\epsilon>0$ of the first-slot flattening of $T_l(Y_0)$.  Then
\[
 \E\rank P\le\epsilon^{-2}\|T_l\|_2^2
 =(1+o(1))\epsilon^{-2}I.
\]
For every $\delta>0$, let $S_\delta$ be the singular-value clipping of
$T_1$ at $\delta$.  Pointwise,
\[
 |P(Y_0)S_\delta(Y_t)|_{\HS}^2
 \le\delta^2\rank P(Y_0).
\]
Stationarity controls the $o(I)$ clipping error for $T_1(Y_t)$.
Hence the normalized limsup of the first expectation in
Equation~\eqref{eq:hierarchy-random-projection} is at most
$\delta^2/\epsilon^2$.  Send $\delta\downarrow0$ to obtain the tail
claim at threshold $\epsilon$ for $T_l$.  No independence between
$P(Y_0)$ and $T_1(Y_t)$ was required.  Tensor symmetry supplies every
other one-slot flattening.
\end{proof}

\subsection{The curvature estimate and its domains}

It remains to prove the curvature assertion of the hierarchy.
We first justify the differentiated operator identity against a test
whose product with the curvature is in $L^2$. The preceding tail
estimate will supply such a test in the strict branch; the gap branch
will use $T_{l+1}$ directly.

\begin{samepage}
\begin{lemma}[Weighted divergence and the differentiated operator]
\label{lem:hierarchy-divergence}
Let $\eta$ be a smooth positive probability density, let
$A=-\Delta-D\log\eta\cdot D$, and put $H=D^2(-\log\eta)$.
Let $P=(P_{i\alpha})$ have one vector slot $i$ and a finite-dimensional
coefficient index $\alpha$.  Suppose every coefficient belongs to
$\Dom A$ and $HP\in L^2(\eta)$.  Then the weighted adjoint divergence
\[
 (D^*P)_\alpha=-\sum_i\partial_iP_{i\alpha}
                   -\sum_i(\partial_i\log\eta)P_{i\alpha}
\]
belongs to $L^2(\eta)$ and satisfies
\begin{equation}\label{eq:hierarchy-divergence-bound}
 \|D^*P\|_2^2\le\|DP\|_2^2+\langle P,HP\rangle_2.
\end{equation}
For every coefficient field $F\in\Dom A$,
\begin{equation}\label{eq:hierarchy-differentiated-operator}
 \langle AF,D^*P\rangle_2
 =\langle DF,AP\rangle_2+\langle DF,HP\rangle_2.
\end{equation}
In particular this identity requires $HP\in L^2$, not $HDF\in L^2$.
\end{lemma}
\end{samepage}

\begin{proof}
For compactly supported smooth $P$, commuting two first derivatives in
weighted integration by parts gives
\begin{equation}\label{eq:hierarchy-divergence-identity}
 \|D^*P\|_2^2
 =\sum_{i,j,\alpha}
       \langle\partial_jP_{i\alpha},\partial_iP_{j\alpha}\rangle_2
   +\langle P,HP\rangle_2.
\end{equation}
Cauchy--Schwarz applied to the two arrays with $i,j$ interchanged bounds
the sum by $\|DP\|_2^2$, without a dimension factor.

For the stated $P$, membership in $\Dom A$ gives $P\in H^2_{\rm loc}$
by local elliptic regularity and $DP\in L^2(\eta)$.  Choose smooth
cutoffs $0\le\chi_R\le1$ tending to one with
$\|D\chi_R\|_\infty\to0$.  Equation~\eqref{eq:hierarchy-divergence-identity}
applies to $\chi_RP$ by local Sobolev approximation.  Since
\[
 \|D(\chi_RP)-DP\|_2\longrightarrow0,
 \qquad
 \langle\chi_RP,H\chi_RP\rangle_2
   \longrightarrow\langle P,HP\rangle_2,
\]
its right side is bounded.  The second convergence uses
$|\langle P,HP\rangle|\in L^1$, a consequence of $P,HP\in L^2$.
Weak compactness in $L^2$ and distributional convergence of
$D^*(\chi_RP)$ identify the limit as $D^*P$.  Lower semicontinuity
proves Equation~\eqref{eq:hierarchy-divergence-bound}.  Only first
derivatives of the cutoffs have entered the estimate.

For compactly supported smooth $F$, differentiate the scalar operator:
\[
 D(AF)=A(DF)+HDF.
\]
Pair this identity with $P$, use symmetry of $H$, and integrate by
parts to obtain Equation~\eqref{eq:hierarchy-differentiated-operator}.
For general $F\in\Dom A$, use the operator core from
Lemma~\ref{lem:dirichlet-core} to choose compact smooth $F_j$ with
$F_j\to F$ and $AF_j\to AF$ in $L^2$.  This convergence also holds
in the form norm, since
\[
 \|D(F_j-F)\|_2^2
 =\langle F_j-F,A(F_j-F)\rangle_2\longrightarrow0.
\]
All terms in Equation~\eqref{eq:hierarchy-differentiated-operator}
therefore converge: the three fixed factors are $D^*P$, $AP$, and
$HP$, each in $L^2$.  This proves the claimed global identity.
\end{proof}

\begin{proof}[Completion of the proof of Proposition~\ref{prop:hierarchy}]
Lemma~\ref{lem:hierarchy-construction} supplies the fields and all
assertions except Equation~\eqref{eq:hierarchy-curvature}.  Fix $l\ge0$
and write $F=T_l$.  View $DF$ as a matrix with its derivative slot as
row index and all old tensor slots as column index.  We construct a
matrix field $P$ satisfying
\begin{equation}\label{eq:hierarchy-curvature-test}
 \|P-DF\|_2+\|(A-1)P\|_2=o(\sqrt I),\qquad
 \|DP\|_2=O(\sqrt I),\qquad \|QP\|_2=o(\sqrt I),
\end{equation}
where $Q=Dw$ in the strict branch and $Q=0$ in the gap branch.

In the strict branch, apply Lemma~\ref{lem:hierarchy-operator-tails}
to the corresponding flattening of $T_{l+1}$.  Choose a singular-value
cutoff $\delta\downarrow0$ sufficiently slowly that the clipped field
$C$ obeys
\[
 \|C-T_{l+1}\|_2=o(\sqrt I),\qquad
 \|C(y)\|_{\op}\le\delta.
\]
Set $P=e^{1-A}C=eP_1C$.  Markov averaging and convexity of the
operator-norm ball give $\|P(y)\|_{\op}\le e\delta$.  The spectral
multipliers $e^{1-x}$, $(x-1)e^{1-x}$, and $\sqrt x\,e^{1-x}$ are
bounded for $x\ge0$.  Moreover
$|e^{1-x}-1|\le C|x-1|$ there.  These bounds, the approximation of
$C$, and the near-eigenfunction assertion for $T_{l+1}$ give the first
two estimates in Equation~\eqref{eq:hierarchy-curvature-test}.
Finally $\|Q\|_2^2=(1+o(1))I$ and
\[
 \|QP\|_2\le\|P\|_{\op,\infty}\|Q\|_2=o(\sqrt I).
\]
The clipping is only used to construct $P$, so it need not preserve
tensor symmetry.  In the gap branch take $P=T_{l+1}$ directly.
The first two estimates follow from
Equation~\eqref{eq:hierarchy-eigen}, and the last is immediate from
$Q=0$.  No operator-tail argument is needed in this branch.

The drift hypothesis gives $H=D^2(-\log\eta)=K-Q$ in either case.
Both $KP$ and $QP$ are in $L^2$, and the operator-domain assertion for
$P$ was retained by its construction.  Lemma~\ref{lem:hierarchy-divergence}
therefore applies.  In particular,
\[
 \|D^*P\|_2^2
 \le\|DP\|_2^2+s^{-1}\|P\|_2^2+\|P\|_2\|QP\|_2=O(I).
\]
Insert $AF=F+(A-1)F$ and $AP=P+(A-1)P$ into
Equation~\eqref{eq:hierarchy-differentiated-operator}.  The terms
$\langle F,D^*P\rangle_2$ and $\langle DF,P\rangle_2$ cancel.
Consequently
\begin{align*}
 \langle DF,KP\rangle_2
 ={}&\langle(A-1)F,D^*P\rangle_2
     -\langle DF,(A-1)P\rangle_2
     +\langle DF,QP\rangle_2
 =o(I).
\end{align*}
Since $\|K\|_{\op,\infty}\le s^{-1}$ and
$P-DF=o_{L^2}(\sqrt I)$, this gives
\[
 \langle DF,KDF\rangle_2=o(I).
\]
Replacing $DF$ by $T_{l+1}$ costs $o(I)$ by the derivative estimate.
It proves Equation~\eqref{eq:hierarchy-curvature} at level $l+1$ on
the derivative slot, and hence on every slot by symmetry.  Level zero
was an initial assumption.  All claims of the proposition follow.
\end{proof}

To use many tensor degrees at once, first impose the estimates on a
fixed finite list of levels, projection-rank constants, and times.
Then increase that list sufficiently slowly along the original sequence.
For example, one may choose $k\to\infty$ so that the normalized errors
in Proposition~\ref{prop:hierarchy} tend to zero uniformly through
level $3k$, while $\sup_{l\le3k}\|T_l\|_2^2/I$ stays bounded.
The same diagonal applies to any fixed number of terms in
Lemma~\ref{lem:hermite-expansion} and any fixed bounded time interval.
Thus averaging over a band of degrees or shifting that band by a fixed
integer does not change any of the preceding normalized estimates.
